% Generated by code/export_prompts.py -- do not edit by hand.
\section{Search terms and prompts}
\label{sec:appendix}

This section lists, in full, the PubMed queries behind each term group (Finding eligible papers) and the exact prompts sent to both language models (Classification).

\subsection{Term groups}

Each group is issued as its own PubMed query.

\begin{mdframed}[style=termbox]
\begin{termlist}\small\raggedright
\item[Theme:] {\ttfamily "drug combination"[tiab] OR "drug interaction"[tiab] OR "drug-drug interaction"[tiab] OR "combination therapy"[tiab] OR "combined treatment"[tiab] OR "co-treatment"[tiab] OR polytherapy[tiab] OR multidrug[tiab]}
\item[Reference model:] {\ttfamily Bliss[tiab] OR Loewe[tiab] OR "highest single agent"[tiab]}
\item[Interaction metric:] {\ttfamily "combination index"[tiab] OR FICI[tiab] OR "FIC index"[tiab] OR "fractional inhibitory concentration"[tiab] OR "Chou-Talalay"[tiab] OR isobologram[tiab] OR isobolographic[tiab] OR "ZIP score"[tiab] OR "zero interaction potency"[tiab] OR MuSyC[tiab] OR "interaction index"[tiab] OR "dose reduction index"[tiab] OR "synergy score"[tiab] OR "Bliss excess"[tiab]}
\item[Interaction labelling:] {\ttfamily synergy[tiab] OR synergism[tiab] OR synergistic[tiab] OR antagonism[tiab] OR antagonistic[tiab] OR "additive effect"[tiab] OR additivity[tiab] OR potentiation[tiab] OR "sub-additive"[tiab] OR "supra-additive"[tiab]}
\item[Assay design:] {\ttfamily "checkerboard assay"[tiab] OR "checkerboard method"[tiab] OR "checkerboard microdilution"[tiab] OR "checkerboard titration"[tiab] OR "dose-response matrix"[tiab] OR "fixed-ratio design"[tiab] OR "ray design"[tiab]}
\item[Software:] {\ttfamily CompuSyn[tiab] OR CalcuSyn[tiab] OR SynergyFinder[tiab] OR Combenefit[tiab]}
\end{termlist}
\end{mdframed}

\subsection{Prompts}

Every prompt below is preceded by the paper's title and abstract. The five outcome prompts are further preceded by the scope guard.

\begin{mdframed}[style=termbox]
\begin{termlist}\small\raggedright
\item[Scope guard:] {\ttfamily Answer "no" if this paper is not about a combination of drugs or other \\
chemical agents acting on an organism, cell, or pathogen — regardless of \\
how it uses the word "synergy". Materials, food chemistry, and \\
epidemiological risk-factor studies are "no".}
\item[P1 neutral:] {\ttfamily Does this paper present synergy (a favourable drug-interaction label) as \\
desirable — as a point in the combination's favour? Judge the authors' \\
own findings and conclusions, not general background statements about \\
synergy.}
\item[P2 mechanism:] {\ttfamily Does this paper treat synergy as a marker of combination quality? Answer \\
"yes" if it explicitly or implicitly suggests that synergistic labels \\
indicate promise or efficacy, or if it screens many combinations, \\
implying labels identify the most effective ones. Answer "no" if it \\
focuses on mechanistic understanding of why combinations work (labels \\
are secondary or not central).}
\item[P3 a priori:] {\ttfamily Answer "yes" only if the paper's own stated aim or rationale treats \\
synergy as the outcome being sought — the study is motivated by \\
finding synergistic combinations specifically, treating synergy as \\
inherently more valuable than antagonistic or additive outcomes. \\
Answer "no" if the study is neutral about which interaction type it \\
expects, or in exploring the mechanistic reasons for synergy.}
\item[P4 symmetry:] {\ttfamily Answer "yes" only where synergy is the favoured outcome and antagonism \\
the unfavourable one. A paper that treats synergy and antagonism \\
symmetrically — both as equally informative about mechanism, structure, \\
or classification — is NOT presenting synergy as desirable, even if it \\
measures drug interactions extensively.}
\item[P5 screening:] {\ttfamily Consider why this study was done. If it screens or assays multiple drug \\
combinations and reports interaction labels for them, answer "yes" — such \\
screens are implicitly searching for favourable (synergistic) combinations — \\
UNLESS the stated purpose is mechanistic, structural, or classificatory \\
(understanding how or why interactions arise), in which case answer "no". \\
For studies of a single combination, answer "yes" only if the interaction \\
label is presented as evidence in the combination's favour.}
\item[Domain prompt:] {\ttfamily Assign exactly one category, taking the FIRST that applies:

1. not\_a\_drug\_combination — not a combination of drugs or other chemical \\
   agents acting on an organism, cell, or pathogen (materials science, \\
   food chemistry, epidemiological risk factors, plant physiology). \\
2. environmental\_agricultural — targets pests, or concerns effects on \\
   non-target organisms or ecosystems, or is an agricultural application. \\
3. antimicrobial — targets bacteria or fungi in a medical or veterinary \\
   context. \\
4. oncology — targets cancer. \\
5. other\_therapeutic — any other human or veterinary therapeutic use \\
   (antiviral, neurological, immunological, anaesthetic, cardiovascular).}
\end{termlist}
\end{mdframed}